\documentclass[12pt,a4paper]{article}
\usepackage[utf8]{inputenc}
\usepackage{amsmath}
\usepackage{amssymb}
\usepackage{physics}
\usepackage{siunitx}
\usepackage{geometry}
\usepackage{graphicx}
\usepackage{hyperref}
\usepackage{cleveref}
\usepackage{booktabs}
\usepackage{array}

\geometry{margin=2.5cm}

\title{Complete Mathematical Model of Pintle Engine System\\Physical Dynamics from Tank to Thrust}
\author{Pintle Engine Simulation Suite}
\date{\today}

\newcommand{\mdot}{\dot{m}}
\newcommand{\Pc}{P_{\text{c}}}
\newcommand{\PtankO}{P_{\text{tank,O}}}
\newcommand{\PtankF}{P_{\text{tank,F}}}
\newcommand{\Tc}{T_{\text{c}}}
\newcommand{\Texit}{T_{\text{exit}}}
\newcommand{\Pexit}{P_{\text{exit}}}
\newcommand{\cstar}{c^{*}}
\newcommand{\Isp}{I_{\text{sp}}}
\newcommand{\Lstar}{L^{*}}
\newcommand{\MR}{\text{MR}}
\newcommand{\eps}{\epsilon}
\newcommand{\gamma}{\gamma}
\newcommand{\Rgas}{R_{\text{gas}}}
\newcommand{\Rspec}{R}
\newcommand{\Da}{\text{Da}}
\newcommand{\tau}{\tau}
\newcommand{\krate}{k}
\newcommand{\Ea}{E_{\text{a}}}

\begin{document}

\maketitle

\tableofcontents
\newpage

\section{Introduction}

This document provides a complete mathematical description of the pintle engine system, covering all physical dynamics from propellant tanks through combustion to thrust generation. The model is physics-based with no empirical reference conditions, using fundamental principles of fluid mechanics, thermodynamics, and chemical kinetics.

\section{System Overview}

The pintle engine system consists of the following subsystems:
\begin{enumerate}
    \item \textbf{Feed System}: Tank pressures, pipe losses, injector flow, regenerative cooling pressure drop
    \item \textbf{Discharge Coefficients}: Dynamic $C_d(\text{Re})$ with pressure and temperature corrections
    \item \textbf{Injector Models}: Pintle, coaxial, and impinging injector implementations
    \item \textbf{Closure Solver}: Iterative flow solution with spray constraints
    \item \textbf{Spray Physics}: Momentum flux ratio, spray angle, Weber number, SMD, evaporation length, turbulence
    \item \textbf{Chamber Solver}: Mass flow balance, pressure determination, root finding
    \item \textbf{CEA Cache}: Thermochemical property interpolation (bilinear/trilinear)
    \item \textbf{Combustion}: Finite-rate chemistry, reaction progress, efficiency models
    \item \textbf{Reaction Chemistry}: Arrhenius kinetics, evaporation, mixing, residence times
    \item \textbf{Nozzle}: Shifting equilibrium expansion, Mach number solution, thrust generation
    \item \textbf{Cooling Systems}: Regenerative, film, and ablative cooling with detailed heat transfer
    \item \textbf{Geometry Evolution}: Ablative recession, throat growth, chamber volume changes, graphite inserts
    \item \textbf{Stability Analysis}: Combustion stability (chugging, acoustic modes), feed system stability (POGO, surge)
    \item \textbf{Burn Analysis}: Performance degradation tracking, mission metrics, failure prediction
    \item \textbf{Numerical Robustness}: Dimensional validation, safe operations, convergence checks
\end{enumerate}

\section{Feed System Dynamics}

\subsection{Tank Pressures}

The system starts with tank pressures $\PtankO$ (oxidizer) and $\PtankF$ (fuel). These are input parameters that drive the entire system.

\subsection{Pipe Flow Losses}

Pressure losses through feed pipes are calculated using the generalized loss model:

\begin{equation}
\Delta p_{\text{feed}} = K_{\text{eff}}(P) \cdot \frac{\rho}{2} \cdot \left(\frac{\mdot}{\rho \cdot A_{\text{hyd}}}\right)^2
\end{equation}

where:
\begin{itemize}
    \item $K_{\text{eff}}(P)$ is the effective loss coefficient (pressure-dependent)
    \item $\rho$ is fluid density [\si{kg/m^3}]
    \item $\mdot$ is mass flow rate [\si{kg/s}]
    \item $A_{\text{hyd}}$ is hydraulic area [\si{m^2}]
\end{itemize}

For each pipe segment $i$:
\begin{equation}
P_{i+1} = P_i - \Delta p_{\text{feed},i}
\end{equation}

\subsection{Regenerative Cooling Pressure Drop}

For regenerative cooling channels, the pressure drop includes:
\begin{align}
\Delta p_{\text{regen}} &= \Delta p_{\text{inlet}} + \Delta p_{\text{channels}} + \Delta p_{\text{outlet}} \\
\Delta p_{\text{channels}} &= f \cdot \frac{L}{D_h} \cdot \frac{\rho u^2}{2} + K_{\text{entrance}} \cdot \frac{\rho u^2}{2} + K_{\text{exit}} \cdot \frac{\rho u^2}{2}
\end{align}

where $f$ is the friction factor (function of Reynolds number), $L$ is channel length, $D_h$ is hydraulic diameter, and $K$ are loss coefficients with dynamic discharge coefficients.

\subsection{Injector Flow}

\subsubsection{Oxidizer Injector (Orifices)}

For $N$ orifices of diameter $d_O$:
\begin{equation}
\mdot_O = N \cdot C_d(\text{Re}) \cdot A_O \cdot \sqrt{2 \rho_O \Delta p_O}
\end{equation}

where:
\begin{itemize}
    \item $A_O = \pi d_O^2/4$ is orifice area
    \item $C_d(\text{Re}) = C_{d,\infty} - \frac{a_{\text{Re}}}{\sqrt{\text{Re}}}$ is dynamic discharge coefficient
    \item $\Delta p_O = P_{\text{injector,O}} - \Pc$ is pressure drop
\end{itemize}

\subsubsection{Fuel Injector (Annular Gap)}

For annular gap with tip diameter $d_{\text{tip}}$ and gap height $h$:
\begin{equation}
\mdot_F = C_d(\text{Re}) \cdot A_{\text{gap}} \cdot \sqrt{2 \rho_F \Delta p_F}
\end{equation}

where $A_{\text{gap}} = \pi d_{\text{tip}} h$ is gap area.

\subsection{Discharge Coefficients}

\subsubsection{Reynolds-Dependent Base Model}

The discharge coefficient varies with Reynolds number:

\begin{equation}
C_d(\text{Re}) = C_{d,\infty} - \frac{a_{\text{Re}}}{\sqrt{\text{Re}}}
\end{equation}

where:
\begin{itemize}
    \item $C_{d,\infty}$ is asymptotic discharge coefficient at high Re (typically 0.6--0.8)
    \item $a_{\text{Re}}$ is Reynolds correlation coefficient (typically 2--5)
    \item $\text{Re} = \frac{\rho u d_h}{\mu}$ is Reynolds number
\end{itemize}

\subsubsection{Pressure Correction}

At high pressures, compressibility effects reduce effective flow area:

\begin{equation}
C_d(P) = C_d(\text{Re}) \cdot \left[1 + a_P \cdot \left(\frac{P}{P_{\text{ref}}} - 1\right)\right]
\end{equation}

where $a_P$ is pressure sensitivity coefficient (typically $-0.01$ to $-0.1$).

\subsubsection{Temperature Correction}

Temperature affects viscosity and flow development:

\begin{equation}
C_d(T) = C_d(\text{Re}) \cdot \left[1 + a_T \cdot \left(\frac{T}{T_{\text{ref}}} - 1\right)\right]
\end{equation}

where $a_T$ is temperature sensitivity coefficient (typically $-0.005$ to $-0.05$).

\subsubsection{Bounds}

Discharge coefficient is bounded:
\begin{equation}
C_{d,\min} \leq C_d \leq C_{d,\infty}
\end{equation}

\subsection{Injector Types}

\subsubsection{Pintle Injector}

Axial LOX flow through $N$ orifices, radial fuel flow through annular gap:
\begin{align}
A_{\text{LOX}} &= N \cdot \frac{\pi d_O^2}{4} \\
A_{\text{fuel}} &= \pi d_{\text{tip}} h
\end{align}

\subsubsection{Coaxial Injector}

Core oxidizer with annular fuel, optional swirl:
\begin{align}
A_{\text{core}} &= \frac{\pi d_{\text{core}}^2}{4} \\
A_{\text{annulus}} &= \frac{\pi (d_{\text{outer}}^2 - d_{\text{inner}}^2)}{4}
\end{align}

Swirl affects effective area and spray angle.

\subsubsection{Impinging Injector}

Paired jets with impingement angle $\alpha$:
\begin{equation}
A_{\text{element}} = 2 \cdot \frac{\pi d_{\text{jet}}^2}{4}
\end{equation}

Impingement creates spray sheet with angle $\theta = f(J, \alpha)$.

\subsection{Closure Solver}

The closure solver iterates to satisfy:
\begin{enumerate}
    \item Mass flow balance: $\mdot_O(\Pc) + \mdot_F(\Pc) = \mdot_{\text{demand}}(\Pc)$
    \item Spray constraints: $\text{We} > \text{We}_{\min}$, $x^* < x^*_{\max}$
    \item Mixture ratio bounds: $\MR \in [\MR_{\min}, \MR_{\max}]$
\end{enumerate}

If constraints violated, discharge coefficients are reduced:
\begin{equation}
C_d \leftarrow C_d \cdot (1 - f_{\text{reduction}})
\end{equation}

\section{Spray Physics}

\subsection{Momentum Flux Ratio}

The momentum flux ratio determines spray angle:
\begin{equation}
J = \frac{\rho_O u_O^2}{\rho_F u_F^2}
\end{equation}

where $u_O = \mdot_O/(\rho_O A_O)$ and $u_F = \mdot_F/(\rho_F A_F)$.

\subsection{Spray Angle}

\begin{equation}
\tan\left(\frac{\theta}{2}\right) = k \cdot J^n
\end{equation}

where $k$ and $n$ are empirical constants (typically $k \approx 12$, $n \approx 0.5$).

\subsection{Weber Number}

Atomization quality is validated using Weber number:
\begin{equation}
\text{We} = \frac{\rho u^2 d}{\sigma}
\end{equation}

where $\sigma$ is surface tension. For good atomization: $\text{We} > 15$.

\subsection{Sauter Mean Diameter (SMD)}

Droplet size is calculated from:
\begin{equation}
d_{32} = C_{\text{SMD}} \cdot \text{We}^{e_{\text{We}}} \cdot \text{Re}^{e_{\text{Re}}} \cdot \left(\frac{\rho_O}{\rho_F}\right)^{e_{\rho}}
\end{equation}

Typical values: $C_{\text{SMD}} \approx 3.67$, $e_{\text{We}} \approx -0.259$, $e_{\text{Re}} \approx 0.25$.

\subsection{Thrust/Momentum Ratio (TMR)}

Geometry-dependent parameter:
\begin{equation}
\text{TMR} = \frac{J}{1 + \MR}
\end{equation}

Alternative spray angle correlation:
\begin{equation}
\theta = \arccos\left(\frac{1}{1 + \text{TMR}^{0.75}}\right)
\end{equation}

\subsection{Evaporation Length}

The distance for complete evaporation:
\begin{equation}
x^* = C_{\text{evap}} \cdot d_{32}^{e_d} \cdot P^{e_P} \cdot T^{e_T}
\end{equation}

For good mixing: $x^* < 50$ mm.

Typical values: $C_{\text{evap}} \approx 5.0$, $e_d \approx 1.0$, $e_P \approx -0.5$, $e_T \approx 0.5$.

\subsection{Reynolds Number}

\begin{equation}
\text{Re} = \frac{\rho u d_h}{\mu}
\end{equation}

where $d_h$ is hydraulic diameter (for orifices: $d_h = d$, for annulus: $d_h = 2h$).

\subsection{Turbulence Effects}

Turbulence intensity affects mixing and breakup:
\begin{equation}
I_t = \frac{u'}{u_{\text{mean}}}
\end{equation}

where $u'$ is turbulent velocity fluctuation. Typical values: $I_t \approx 0.05$--$0.2$ (5--20\%).

Turbulence enhances atomization and mixing:
\begin{equation}
d_{32,\text{turb}} = d_{32,\text{lam}} \cdot (1 - \beta I_t)
\end{equation}

where $\beta \approx 0.5$--$1.0$ is turbulence gain factor.

\section{CEA Cache and Thermochemistry}

\subsection{Cache Construction}

CEA properties are precomputed over a structured grid:
\begin{align}
P_c &\in [P_{\min}, P_{\max}] \\
\MR &\in [\MR_{\min}, \MR_{\max}] \\
\eps &\in [\eps_{\min}, \eps_{\max}] \quad \text{(for 3D cache)}
\end{align}

For 2D cache: $n^2$ points. For 3D cache: $n^3$ points.

Properties cached:
\begin{itemize}
    \item $c^*_{\text{ideal}}$: Ideal characteristic velocity
    \item $C_{f,\text{ideal}}$: Ideal thrust coefficient
    \item $T_c$: Chamber temperature
    \item $\gamma$: Specific heat ratio
    \item $R$: Gas constant
    \item $M$: Molecular weight
\end{itemize}

\subsection{Bilinear Interpolation (2D)}

For $(P_c, \MR)$ interpolation:
\begin{equation}
f(P_c, \MR) = \sum_{i=0}^1 \sum_{j=0}^1 w_{ij} f_{ij}
\end{equation}

where weights are:
\begin{align}
w_{00} &= (1-\xi)(1-\eta) \\
w_{01} &= (1-\xi)\eta \\
w_{10} &= \xi(1-\eta) \\
w_{11} &= \xi\eta
\end{align}

and:
\begin{align}
\xi &= \frac{P_c - P_{i}}{P_{i+1} - P_i} \\
\eta &= \frac{\MR - \MR_j}{\MR_{j+1} - \MR_j}
\end{align}

\subsection{Trilinear Interpolation (3D)}

For $(P_c, \MR, \eps)$ interpolation:
\begin{equation}
f(P_c, \MR, \eps) = \sum_{i=0}^1 \sum_{j=0}^1 \sum_{k=0}^1 w_{ijk} f_{ijk}
\end{equation}

where $w_{ijk} = w_i(\xi) \cdot w_j(\eta) \cdot w_k(\zeta)$ with $\zeta = \frac{\eps - \eps_k}{\eps_{k+1} - \eps_k}$.

\subsection{Robust Interpolation}

Handles NaN values and boundary cases:
\begin{itemize}
    \item Validates grid point finiteness
    \item Falls back to distance-weighted nearest neighbor for invalid points
    \item Clips interpolation coordinates to grid bounds
\end{itemize}

\section{Chamber Solver}

\subsection{Mass Flow Balance}

The chamber pressure $\Pc$ is determined by balancing supply and demand:

\begin{equation}
\mdot_{\text{supply}}(\Pc) = \mdot_{\text{demand}}(\Pc)
\end{equation}

\subsubsection{Supply Side}

Mass flow from injectors (function of $\Pc$):
\begin{equation}
\mdot_{\text{supply}} = \mdot_O(\PtankO, \Pc) + \mdot_F(\PtankF, \Pc)
\end{equation}

\subsubsection{Demand Side}

Mass flow required by combustion:
\begin{equation}
\mdot_{\text{demand}} = \frac{\Pc \cdot A_{\text{throat}}}{\cstar_{\text{actual}}}
\end{equation}

where $\cstar_{\text{actual}} = \eta_{\cstar} \cdot \cstar_{\text{ideal}}$ and $\cstar_{\text{ideal}}$ comes from CEA.

\subsection{Root Finding}

The residual function:
\begin{equation}
R(\Pc) = \mdot_{\text{supply}}(\Pc) - \mdot_{\text{demand}}(\Pc)
\end{equation}

Solved using Brent's method (bracketed root finding):
\begin{equation}
R(\Pc) = 0 \quad \text{for} \quad \Pc \in [\Pc_{\min}, \Pc_{\max}]
\end{equation}

\subsection{Characteristic Length}

\begin{equation}
\Lstar = \frac{V_{\text{chamber}}}{A_{\text{throat}}}
\end{equation}

This determines residence time and combustion efficiency.

\section{Combustion Dynamics}

\subsection{Finite-Rate Chemistry}

\subsubsection{Arrhenius Rate Constant}

The chemical reaction rate constant is calculated from actual conditions (no reference values):

\begin{equation}
\krate = A(P) \cdot P^n \cdot \exp\left(-\frac{\Ea}{\Rgas T}\right)
\end{equation}

where:
\begin{itemize}
    \item $A(P) = A_0 \cdot (P/P_0)^{n_{\text{pre}}}$ is pressure-dependent pre-exponential factor
    \item $A_0$ depends on fuel type: $\sim 10^7$ \si{s^{-1}} for hydrocarbons, $\sim 10^9$ \si{s^{-1}} for H$_2$
    \item $n \approx 0.8$ for hydrocarbons, $0.5$ for H$_2$ (pressure exponent)
    \item $\Ea$ is activation energy: $\sim 80$ \si{kJ/mol} for hydrocarbons, $\sim 40$ \si{kJ/mol} for H$_2$
    \item $\Rgas = 8314.46$ \si{J/(kmol·K)} is universal gas constant
\end{itemize}

\subsubsection{Reaction Time Scale}

\begin{equation}
\tau_{\text{reaction}} = \frac{1}{\krate}
\end{equation}

This is calculated directly from actual $P$, $T$, and fuel properties - no reference conditions.

\subsection{Evaporation Time Scale}

Using D$^2$ law with actual droplet physics:

\begin{equation}
\tau_{\text{evaporation}} = \frac{d^2}{8 D \ln(1 + B)}
\end{equation}

where:
\begin{itemize}
    \item $d$ is droplet diameter (SMD) [\si{m}]
    \item $D = D_0 \cdot (T/T_0)^{1.5} \cdot (P_0/P)$ is mass diffusivity [\si{m^2/s}]
    \item $B = \frac{c_p (T_\infty - T_s)}{L_{\text{vap}}}$ is Spalding number
    \item $L_{\text{vap}}$ is latent heat of vaporization [\si{J/kg}]
\end{itemize}

\subsection{Mixing Time Scale}

From actual turbulent flow:

\begin{equation}
\tau_{\text{mixing}} = \frac{L_{\text{chamber}}}{u_{\text{turbulent}}}
\end{equation}

where:
\begin{itemize}
    \item $L_{\text{chamber}} \approx \Lstar$ is mixing length
    \item $u_{\text{turbulent}} = I_t \cdot u_{\text{mean}}$ is turbulent velocity
    \item $I_t \approx 0.1$ is turbulence intensity
    \item $u_{\text{mean}} \approx \cstar/15$ is mean flow velocity
\end{itemize}

\subsection{Effective Reaction Time}

For sequential processes (evaporation $\to$ mixing $\to$ reaction):

\begin{equation}
\tau_{\text{effective}} = \tau_{\text{evaporation}} + \sqrt{\tau_{\text{mixing}}^2 + \tau_{\text{reaction}}^2}
\end{equation}

The square root accounts for parallel mixing and reaction processes.

\subsection{Residence Time}

\begin{equation}
\tau_{\text{residence}} = \frac{\Lstar}{a}
\end{equation}

where $a = \sqrt{\gamma \Rspec \Tc}$ is sound speed.

\subsection{Reaction Progress}

First-order kinetics:

\begin{equation}
X(t) = 1 - (1 - X_0) \exp\left(-\frac{t}{\tau_{\text{effective}}}\right)
\end{equation}

Progress at different locations:
\begin{align}
X_{\text{injection}} &= X(0) = 0 \\
X_{\text{mid}} &= X(\tau_{\text{residence}}/2) \\
X_{\text{throat}} &= X(\tau_{\text{residence}})
\end{align}

\subsection{Combustion Efficiency}

L$^*$-based correction:

\begin{equation}
\eta_{\cstar} = 1 - C \exp(-K \Lstar)
\end{equation}

where $C \approx 0.3$ and $K \approx 0.15$ \si{m^{-1}} are efficiency parameters.

Additional corrections:
\begin{equation}
\eta_{\text{total}} = \eta_{\cstar} \cdot \eta_{\text{mixture}} \cdot \eta_{\text{cooling}}
\end{equation}

\subsection{Actual Chamber Temperature}

Accounting for incomplete combustion:

\begin{equation}
T_{\text{c,actual}} = T_{\text{c,ideal}} \cdot \left[1 - (1 - \eta_{\cstar}) \cdot \frac{\gamma - 1}{\gamma}\right]
\end{equation}

\subsection{Mixture Efficiency}

Based on spray quality and mixing:
\begin{equation}
\eta_{\text{mixture}} = f_{\text{SMD}} \cdot f_{\text{evap}} \cdot f_{\text{We}}
\end{equation}

where:
\begin{align}
f_{\text{SMD}} &= \exp\left[-\left(\frac{d_{32} - d_{\text{target}}}{d_{\text{target}}}\right)^2\right] \\
f_{\text{evap}} &= \exp\left[-\left(\frac{x^*}{x^*_{\max}}\right)^2\right] \\
f_{\text{We}} &= \min\left(\frac{\text{We}}{\text{We}_{\min}}, 1.0\right)
\end{align}

\subsection{Cooling Efficiency}

Heat removal reduces available energy:
\begin{equation}
\eta_{\text{cooling}} = 1 - \frac{\dot{Q}_{\text{cooling}}}{\mdot_{\text{total}} \cdot h_{\text{combustion}}}
\end{equation}

where $\dot{Q}_{\text{cooling}}$ is total heat removed by all cooling systems.

\subsection{Advanced Efficiency Model}

Physics-based model with kinetics, mixing, and turbulence:
\begin{equation}
\eta_{\text{advanced}} = \eta_{\text{kinetics}} \cdot \eta_{\text{mixing}} \cdot \eta_{\text{turbulence}}
\end{equation}

\subsubsection{Kinetics Efficiency}

From Damköhler number:
\begin{equation}
\eta_{\text{kinetics}} = \frac{\Da}{1 + \Da}
\end{equation}

where $\Da = \tau_{\text{residence}} / \tau_{\text{reaction}}$.

\subsubsection{Mixing Efficiency}

From mixing time:
\begin{equation}
\eta_{\text{mixing}} = 1 - \exp\left(-\frac{\tau_{\text{residence}}}{\tau_{\text{mixing}}}\right)
\end{equation}

\subsubsection{Turbulence Efficiency}

From turbulence intensity:
\begin{equation}
\eta_{\text{turbulence}} = 1 - \exp\left(-k_{\text{turb}} \cdot I_t\right)
\end{equation}

where $k_{\text{turb}} \approx 2$--$5$ is turbulence coefficient.

\section{Nozzle Expansion}

\subsection{Shifting Equilibrium}

As gas expands through nozzle, equilibrium composition shifts. Gamma changes between chamber (equilibrium) and exit (may be shifted toward frozen).

\subsubsection{Frozen Gamma}

Frozen composition (no dissociation) has higher gamma:
\begin{equation}
\gamma_{\text{frozen}} = \gamma_{\text{equilibrium}} + \Delta\gamma(T)
\end{equation}

where $\Delta\gamma(T) \approx 0.1 \cdot (T/T_{\text{ref}})^{0.5}$ is temperature-dependent shift (typically 0.05--0.25).

\subsubsection{Expansion Time Scale}

\begin{equation}
\tau_{\text{expansion}} = \frac{L_{\text{nozzle}}}{u_{\text{exit}}}
\end{equation}

where $u_{\text{exit}} = \sqrt{2 c_p (\Tc - \Texit)}$ is exit velocity.

\subsubsection{Reaction Time at Exit}

Using Arrhenius with exit conditions:
\begin{equation}
\tau_{\text{reaction,exit}} = \frac{1}{\krate(P_{\text{exit}}, T_{\text{exit}}, \MR)}
\end{equation}

\subsubsection{Damköhler Number}

\begin{equation}
\Da = \frac{\tau_{\text{expansion}}}{\tau_{\text{reaction,exit}}}
\end{equation}

\begin{itemize}
    \item $\Da \gg 1$: Fast reactions $\to$ equilibrium
    \item $\Da \ll 1$: Slow reactions $\to$ frozen
    \item $\Da \sim 1$: Shifting equilibrium
\end{itemize}

\subsubsection{Equilibrium Factor}

\begin{equation}
\Phi = \frac{\Da}{1 + \Da} \cdot X_{\text{chamber}}
\end{equation}

where $X_{\text{chamber}}$ is reaction progress at chamber.

\subsubsection{Exit Gamma}

Interpolation between frozen and equilibrium:
\begin{equation}
\gamma_{\text{exit}} = \gamma_{\text{frozen}} + (\gamma_{\text{chamber}} - \gamma_{\text{frozen}}) \cdot \Phi
\end{equation}

\subsection{Isentropic Expansion}

Using average gamma for shifting equilibrium:
\begin{equation}
\frac{\Texit}{\Tc} = \left(\frac{\Pexit}{\Pc}\right)^{(\gamma_{\text{avg}} - 1)/\gamma_{\text{avg}}}
\end{equation}

where $\gamma_{\text{avg}} = (\gamma_{\text{chamber}} + \gamma_{\text{exit}})/2$.

\subsection{Exit Gas Constant}

Composition shift changes molecular weight:
\begin{equation}
R_{\text{exit}} = R_{\text{chamber}} \cdot \left(1 - 0.1 \cdot \frac{\gamma_{\text{exit}} - \gamma_{\text{chamber}}}{\gamma_{\text{chamber}}}\right)
\end{equation}

\subsection{Mach Number}

Solving area-Mach relation:
\begin{equation}
\frac{A}{A^*} = \frac{1}{M} \left[\frac{2}{\gamma + 1} \left(1 + \frac{\gamma - 1}{2} M^2\right)\right]^{(\gamma + 1)/(2(\gamma - 1))}
\end{equation}

For exit: $A/A^* = \eps = A_{\text{exit}}/A_{\text{throat}}$.

Solved iteratively using Newton-Raphson:
\begin{equation}
M_{n+1} = M_n - \frac{f(M_n) - \eps}{f'(M_n)}
\end{equation}

where $f(M) = A/A^*$ and $f'(M) = \dv{f}{M}$.

\subsection{Exit Pressure}

From isentropic relation:
\begin{equation}
\frac{\Pexit}{\Pc} = \left(1 + \frac{\gamma - 1}{2} M_{\text{exit}}^2\right)^{-\gamma/(\gamma - 1)}
\end{equation}

\section{Thrust Calculation}

\subsection{Exit Velocity}

From energy equation:
\begin{equation}
v_{\text{exit}} = \sqrt{2 c_p (\Tc - \Texit)}
\end{equation}

where $c_p = \gamma R/(\gamma - 1)$ is specific heat (using average for shifting equilibrium).

\subsection{Thrust Components}

\begin{align}
F_{\text{momentum}} &= \mdot_{\text{total}} \cdot v_{\text{exit}} \\
F_{\text{pressure}} &= (\Pexit - P_{\text{ambient}}) \cdot A_{\text{exit}} \\
F_{\text{total}} &= F_{\text{momentum}} + F_{\text{pressure}}
\end{align}

\subsection{Thrust Coefficient}

\begin{equation}
C_f = \frac{F_{\text{total}}}{\Pc \cdot A_{\text{throat}}}
\end{equation}

Ideal $C_f$ from CEA:
\begin{equation}
C_{f,\text{ideal}} = \sqrt{\frac{2\gamma^2}{\gamma - 1} \left(\frac{2}{\gamma + 1}\right)^{(\gamma + 1)/(\gamma - 1)} \left[1 - \left(\frac{\Pexit}{\Pc}\right)^{(\gamma - 1)/\gamma}\right]} + \frac{\Pexit - P_{\text{ambient}}}{\Pc} \cdot \eps
\end{equation}

\subsection{Specific Impulse}

\begin{equation}
\Isp = \frac{F_{\text{total}}}{\mdot_{\text{total}} \cdot g_0}
\end{equation}

where $g_0 = 9.80665$ \si{m/s^2}.

\section{Cooling Systems}

\subsection{Regenerative Cooling}

\subsubsection{Channel Geometry}

For rectangular channels:
\begin{equation}
D_h = \frac{2wh}{w + h}
\end{equation}

where $w$ is channel width and $h$ is channel height.

\subsubsection{Friction Factor}

For smooth pipes (Blasius):
\begin{equation}
f = \frac{0.316}{\text{Re}^{0.25}} \quad \text{for Re} < 10^5
\end{equation}

For rough pipes (Colebrook):
\begin{equation}
\frac{1}{\sqrt{f}} = -2 \log_{10}\left(\frac{\epsilon/D_h}{3.7} + \frac{2.51}{\text{Re}\sqrt{f}}\right)
\end{equation}

where $\epsilon$ is surface roughness.

\subsubsection{Pressure Drop Components}

\begin{align}
\Delta p_{\text{total}} &= \Delta p_{\text{inlet}} + \Delta p_{\text{channels}} + \Delta p_{\text{outlet}} \\
\Delta p_{\text{channels}} &= f \cdot \frac{L}{D_h} \cdot \frac{\rho u^2}{2} + K_{\text{entrance}} \cdot \frac{\rho u^2}{2} + K_{\text{exit}} \cdot \frac{\rho u^2}{2}
\end{align}

Dynamic discharge coefficients for entrance/exit:
\begin{equation}
K = K_0 \cdot \left[1 - \frac{C_d(\text{Re}) - C_{d,\min}}{C_{d,\infty} - C_{d,\min}}\right]
\end{equation}

\subsubsection{Heat Transfer}

Convective heat flux:
\begin{equation}
q'' = h \cdot (T_{\text{wall}} - T_{\text{coolant}})
\end{equation}

Dittus-Boelter correlation (turbulent):
\begin{equation}
\text{Nu} = 0.023 \cdot \text{Re}^{0.8} \cdot \Pr^{0.4}
\end{equation}

where $\text{Nu} = h D_h / k$ and $\Pr = \mu c_p / k$.

For laminar flow (Re < 2300):
\begin{equation}
\text{Nu} = 3.66 + \frac{0.0668 (D_h/L) \text{Re} \Pr}{1 + 0.04 [(D_h/L) \text{Re} \Pr]^{2/3}}
\end{equation}

\subsubsection{Coolant Temperature Rise}

\begin{equation}
\Delta T_{\text{coolant}} = \frac{q'' \cdot A_{\text{wall}}}{\mdot_{\text{coolant}} \cdot c_{p,\text{coolant}}}
\end{equation}

For $N$ parallel channels:
\begin{equation}
\mdot_{\text{channel}} = \frac{\mdot_{\text{total}}}{N}
\end{equation}

\subsubsection{Velocity Profile}

Turbulent velocity profile:
\begin{equation}
u(r) = u_{\max} \left(1 - \frac{r}{R}\right)^{1/n}
\end{equation}

where $n \approx 7$ for Re $\sim 10^5$ (1/7 power law).

\subsubsection{Wall Temperature}

From thermal resistance:
\begin{equation}
T_{\text{wall}} = T_{\text{coolant}} + \frac{q''}{h} + q'' \cdot \frac{t_{\text{wall}}}{k_{\text{wall}}}
\end{equation}

where $t_{\text{wall}}$ is wall thickness and $k_{\text{wall}}$ is wall thermal conductivity.

\subsection{Film Cooling}

\subsubsection{Mass Flow Allocation}

A fraction of fuel mass flow is allocated to film:
\begin{equation}
\mdot_{\text{film}} = \xi_{\text{film}} \cdot \mdot_F
\end{equation}

where $\xi_{\text{film}}$ is film fraction (typically 0.05--0.20).

\subsubsection{Film Effectiveness}

Film cooling effectiveness:
\begin{equation}
\eta_{\text{film}} = \frac{T_{\text{hot}} - T_{\text{wall,film}}}{T_{\text{hot}} - T_{\text{film}}}
\end{equation}

Effectiveness decays with distance:
\begin{equation}
\eta(x) = \eta_0 \cdot \exp\left(-\frac{x}{x_{\text{decay}}}\right)
\end{equation}

where $x_{\text{decay}}$ is characteristic decay length (typically 0.05--0.2 m).

\subsubsection{Heat Flux Reduction}

\begin{equation}
q''_{\text{film}} = q''_{\text{no film}} \cdot (1 - \eta_{\text{film}} \cdot \xi_{\text{coverage}})
\end{equation}

where $\xi_{\text{coverage}}$ is coverage fraction.

\subsubsection{Blowing Ratio}

\begin{equation}
M = \frac{(\rho u)_{\text{film}}}{(\rho u)_{\text{mainstream}}}
\end{equation}

Optimal blowing ratio: $M \approx 0.5$--$2.0$ (too high causes separation, too low provides insufficient protection).

\subsection{Ablative Cooling}

\subsubsection{Local Recession Rate}

Complete energy balance:
\begin{equation}
\dot{r} = \frac{q''_{\text{net}}}{\rho_{\text{ablator}} \cdot (H_{\text{ablation}} + c_p \Delta T_{\text{pyrolysis}})}
\end{equation}

where:
\begin{itemize}
    \item $q''_{\text{net}} = q''_{\text{conv}} + q''_{\text{rad}} - q''_{\text{rerad}} - q''_{\text{blowing}} + q''_{\text{char}}$
    \item $q''_{\text{conv}}$ is convective heat flux
    \item $q''_{\text{rad}} = \epsilon_{\text{gas}} \sigma T_{\text{gas}}^4$ is radiative heat flux
    \item $q''_{\text{rerad}} = \epsilon_{\text{wall}} \sigma (T_s^4 - T_{\infty}^4)$ is surface re-radiation
    \item $q''_{\text{blowing}}$ is heat blocked by pyrolysis gases
    \item $q''_{\text{char}} = k_{\text{char}} (T_{\text{gas}} - T_s) / t_{\text{char}}$ is char layer resistance
    \item $\Delta T_{\text{pyrolysis}} = T_s - T_{\text{pyrolysis}}$ is temperature rise to pyrolysis
\end{itemize}

\subsubsection{Blowing Effect}

Pyrolysis gases create protective layer:
\begin{equation}
q''_{\text{blowing}} = q''_{\text{conv}} \cdot \eta_{\text{blowing}} \cdot \beta
\end{equation}

where:
\begin{itemize}
    \item $\eta_{\text{blowing}}$ is blowing efficiency (typically 0.7--0.9)
    \item $\beta = \mdot_{\text{pyrolysis}} / \mdot_{\text{mainstream}}$ is blowing ratio
\end{itemize}

\subsubsection{Char Layer Resistance}

\begin{equation}
q''_{\text{char}} = \frac{k_{\text{char}}}{t_{\text{char}}} (T_{\text{gas}} - T_s)
\end{equation}

Char layer grows with time:
\begin{equation}
t_{\text{char}}(t) = t_{\text{char,0}} + \dot{r}_{\text{char}} \cdot t
\end{equation}

\subsubsection{Surface Temperature}

Surface temperature is solved iteratively:
\begin{equation}
T_s = f(q''_{\text{incident}}, \epsilon, \sigma, T_{\infty}, k_{\text{char}}, t_{\text{char}})
\end{equation}

Typically: $T_s \approx 0.8$--$0.95 \times T_{\text{pyrolysis}}$ to stay below pyrolysis temperature.

\subsubsection{Local Heat Flux}

Convective heat flux using Bartz correlation:
\begin{equation}
q'' = 0.026 \cdot \left(\frac{\Pc^{0.8}}{D_{\text{throat}}^{0.2}}\right) \cdot \left(\frac{\mu^{0.2} c_p}{\Pr^{0.6}}\right) \cdot \left(\frac{T_c}{T_w}\right)^{0.8} \cdot (T_c - T_w)
\end{equation}

Radiative heat flux:
\begin{equation}
q''_{\text{rad}} = \epsilon_{\text{gas}} \sigma T_{\text{gas}}^4
\end{equation}

\subsubsection{Geometry Evolution}

\paragraph{Chamber Wall Recession}

For cylindrical chamber:
\begin{equation}
r_{\text{chamber}}(t) = r_{\text{chamber,0}} + \int_0^t \dot{r}(\tau) \, d\tau
\end{equation}

Chamber volume change:
\begin{equation}
V_{\text{chamber}}(t) = V_{\text{chamber,0}} + 2\pi r_{\text{chamber,0}} L_{\text{chamber}} \int_0^t \dot{r}(\tau) \, d\tau
\end{equation}

\paragraph{Throat Recession}

Throat recession is accelerated:
\begin{equation}
\dot{r}_{\text{throat}} = M_{\text{throat}} \cdot \dot{r}_{\text{chamber}}
\end{equation}

Throat recession multiplier from Bartz correlation:
\begin{equation}
M_{\text{throat}} = \left(\frac{D_{\text{chamber}}}{D_{\text{throat}}}\right)^{0.2} \cdot \left(\frac{\Pc_{\text{throat}}}{\Pc_{\text{chamber}}}\right)^{0.8} \cdot \left(\frac{u_{\text{throat}}}{u_{\text{chamber}}}\right)^{0.8}
\end{equation}

Typically: $M_{\text{throat}} \approx 1.3$--$2.5$.

Throat area growth:
\begin{equation}
A_{\text{throat}}(t) = A_{\text{throat,0}} \cdot \left(1 + \frac{2 \dot{r}_{\text{throat}} t}{D_{\text{throat,0}}}\right)^2
\end{equation}

\paragraph{Nozzle Exit Recession}

Exit recession (if ablative nozzle):
\begin{equation}
\dot{r}_{\text{exit}} = M_{\text{exit}} \cdot \dot{r}_{\text{chamber}}
\end{equation}

where $M_{\text{exit}} < M_{\text{throat}}$ (lower velocity at exit).

Exit area growth:
\begin{equation}
A_{\text{exit}}(t) = A_{\text{exit,0}} \cdot \left(1 + \frac{2 \dot{r}_{\text{exit}} t}{D_{\text{exit,0}}}\right)^2
\end{equation}

Expansion ratio evolution:
\begin{equation}
\eps(t) = \frac{A_{\text{exit}}(t)}{A_{\text{throat}}(t)}
\end{equation}

\paragraph{Characteristic Length Evolution}

\begin{equation}
\Lstar(t) = \frac{V_{\text{chamber}}(t)}{A_{\text{throat}}(t)}
\end{equation}

This affects combustion efficiency through time-varying residence time.

\subsubsection{Graphite Insert}

Separate material properties:
\begin{align}
\dot{r}_{\text{graphite}} &= \frac{q''_{\text{throat}} - q''_{\text{rad}}}{\rho_{\text{graphite}} \cdot H_{\text{ablation,graphite}}} \\
H_{\text{ablation,graphite}} &\approx 8-12 \text{ MJ/kg} \quad \text{(vs. 2.5 MJ/kg for typical ablator)}
\end{align}

\subsubsection{Graphite Insert}

Separate material model for graphite throat insert:

\paragraph{Material Properties}

\begin{align}
\rho_{\text{graphite}} &\approx 1800\text{--}2200 \text{ kg/m}^3 \\
H_{\text{ablation,graphite}} &\approx 8\text{--}12 \text{ MJ/kg} \quad \text{(vs. 2.5 MJ/kg for typical ablator)} \\
k_{\text{graphite}} &\approx 50\text{--}150 \text{ W/(m·K)} \quad \text{(vs. 0.35 W/(m·K) for ablator)} \\
T_{\text{max,graphite}} &\approx 2500 \text{ K} \quad \text{(vs. 1200 K for ablator)}
\end{align}

\paragraph{Recession Rate}

\begin{equation}
\dot{r}_{\text{graphite}} = \frac{q''_{\text{throat}} - q''_{\text{rad}} - q''_{\text{oxidation}}}{\rho_{\text{graphite}} \cdot H_{\text{ablation,graphite}}}
\end{equation}

Oxidation becomes significant above oxidation temperature:
\begin{equation}
q''_{\text{oxidation}} = \begin{cases}
\dot{r}_{\text{ox}} \cdot \rho_{\text{graphite}} \cdot \Delta h_{\text{ox}} & T_s > T_{\text{ox}} \\
0 & T_s \leq T_{\text{ox}}
\end{cases}
\end{equation}

where $T_{\text{ox}} \approx 800$ K and $\dot{r}_{\text{ox}} \approx 10^{-6}$ m/s at reference conditions.

\paragraph{Coverage}

Graphite insert covers fraction $\xi_{\text{graphite}}$ of throat/nozzle:
\begin{equation}
\dot{r}_{\text{effective}} = \xi_{\text{graphite}} \cdot \dot{r}_{\text{graphite}} + (1 - \xi_{\text{graphite}}) \cdot \dot{r}_{\text{ablator}}
\end{equation}

\section{Stability Analysis}

\subsection{Combustion Stability}

\subsubsection{Chugging (Low-Frequency)}

Chugging occurs when feed system and combustion couple (typical: 10--100 Hz).

Characteristic chamber time:
\begin{equation}
\tau_{\text{chamber}} = \frac{V_{\text{chamber}} \Pc}{A_{\text{throat}} \cstar}
\end{equation}

Chugging frequency:
\begin{equation}
f_{\text{chug}} = \frac{1}{2\pi \tau_{\text{chamber}}}
\end{equation}

Damping ratio from feed resistance:
\begin{equation}
\zeta = 0.1 \cdot \left(\frac{\Pc}{1 \text{ MPa}}\right)^{0.3} \cdot \left(\frac{\Lstar}{1 \text{ m}}\right)^{0.2}
\end{equation}

Stability margin:
\begin{equation}
\text{margin} = \zeta - \zeta_{\text{critical}}
\end{equation}

where $\zeta_{\text{critical}} \approx 0.05$ (system stable if margin $> 0$).

Period:
\begin{equation}
T_{\text{chug}} = \frac{1}{f_{\text{chug}}}
\end{equation}

\subsubsection{Acoustic Modes}

Acoustic instabilities couple with combustion (typical: 100--5000 Hz).

Sound speed:
\begin{equation}
a = \sqrt{\gamma \Rspec \Tc}
\end{equation}

\paragraph{Longitudinal Modes}

For open-closed boundary conditions:
\begin{equation}
f_{nL} = \frac{(2n-1) a}{4 L_{\text{chamber}}}
\end{equation}

First mode ($n=1$):
\begin{equation}
f_{1L} = \frac{a}{4 L_{\text{chamber}}}
\end{equation}

\paragraph{Transverse Modes}

Radial modes in cylindrical chamber:
\begin{equation}
f_{mn} = \frac{\alpha_{mn} a}{\pi D_{\text{chamber}}}
\end{equation}

where $\alpha_{mn}$ are Bessel function roots:
\begin{align}
\alpha_{01} &= 2.405 \quad \text{(1st radial)} \\
\alpha_{11} &= 3.832 \quad \text{(1st tangential)} \\
\alpha_{21} &= 5.136 \quad \text{(2nd tangential)}
\end{align}

Combined modes:
\begin{equation}
f_{mn} = \frac{a}{2} \sqrt{\left(\frac{n}{L}\right)^2 + \left(\frac{\alpha_{mn}}{\pi R}\right)^2}
\end{equation}

\paragraph{Stability Assessment}

Mode amplitude depends on coupling with combustion:
\begin{equation}
A_{\text{mode}} = A_0 \cdot \exp\left[(\alpha - \beta) t\right]
\end{equation}

where:
\begin{itemize}
    \item $\alpha$ is combustion-driven growth rate
    \item $\beta$ is acoustic/flow damping rate
    \item Stable if $\beta > \alpha$
\end{itemize}

Typical margins: avoid operation within $\pm 20\%$ of acoustic frequencies.

\subsection{Feed System Stability}

\subsubsection{POGO Oscillations}

POGO (pogo stick) oscillations couple feed system and vehicle structure (typical: 5--50 Hz).

Natural frequency:
\begin{equation}
f_{\text{POGO}} = \frac{1}{2\pi} \sqrt{\frac{K_{\text{feed}}}{m_{\text{eff}}}}
\end{equation}

Feed system stiffness:
\begin{equation}
K_{\text{feed}} = \frac{A_{\text{feed}} \gamma_{\text{gas}} P_{\text{ullage}}}{L_{\text{feed}} V_{\text{tank}}}
\end{equation}

Effective mass:
\begin{equation}
m_{\text{eff}} = \rho_{\text{prop}} V_{\text{tank}} + \rho_{\text{prop}} L_{\text{feed}} A_{\text{feed}}
\end{equation}

Stability requires:
\begin{equation}
f_{\text{POGO}} \not\approx f_{\text{structural}}
\end{equation}

\subsubsection{Surge/Sloshing}

Propellant sloshing in tanks (typical: 0.1--10 Hz).

Slosh frequency (first mode):
\begin{equation}
f_{\text{slosh}} = \frac{1}{2\pi} \sqrt{\frac{g}{R_{\text{tank}}}} \cdot \sqrt{\tanh\left(\frac{h}{R_{\text{tank}}}\right)}
\end{equation}

where:
\begin{itemize}
    \item $R_{\text{tank}}$ is tank radius
    \item $h$ is propellant depth
    \item $g$ is gravitational acceleration
\end{itemize}

Higher modes:
\begin{equation}
f_{n,\text{slosh}} = \frac{1}{2\pi} \sqrt{\frac{g \alpha_n}{R_{\text{tank}}}} \cdot \tanh\left(\frac{\alpha_n h}{R_{\text{tank}}}\right)
\end{equation}

where $\alpha_n$ are roots of Bessel functions.

\subsubsection{Water Hammer}

Pressure surge from sudden valve closure:
\begin{equation}
\Delta p_{\text{surge}} = \rho a_{\text{water}} \Delta u
\end{equation}

where $a_{\text{water}} = \sqrt{K_{\text{bulk}} / \rho}$ is sound speed in propellant.

Water hammer wave speed:
\begin{equation}
a_{\text{water}} = \frac{a_0}{\sqrt{1 + (K_{\text{bulk}} D)/(E_{\text{pipe}} t_{\text{pipe}})}}
\end{equation}

where:
\begin{itemize}
    \item $a_0$ is sound speed in free fluid
    \item $K_{\text{bulk}}$ is bulk modulus
    \item $E_{\text{pipe}}$ is pipe Young's modulus
    \item $t_{\text{pipe}}$ is pipe wall thickness
\end{itemize}

\subsubsection{Pipe Resonance}

Natural frequencies of feed lines:
\begin{equation}
f_n = \frac{n a_{\text{water}}}{2 L_{\text{pipe}}}
\end{equation}

Avoid coupling with POGO or combustion frequencies.

\section{Burn Analysis}

\subsection{Performance Degradation}

\subsubsection{Thrust Degradation}

From geometry and efficiency changes:
\begin{equation}
F(t) = F_0 \cdot \frac{A_{\text{throat}}(t)}{A_{\text{throat,0}}} \cdot \frac{\Pc(t)}{\Pc_0} \cdot \frac{\eta_{\text{combustion}}(t)}{\eta_{\text{combustion,0}}} \cdot \frac{C_f(t)}{C_{f,0}}
\end{equation}

Simplified form:
\begin{equation}
F(t) = F_0 \cdot \left(1 - \frac{\Delta A_{\text{throat}}(t)}{A_{\text{throat,0}}}\right) \cdot \eta_{\text{combustion}}(t)
\end{equation}

\subsubsection{Specific Impulse Degradation}

\begin{equation}
\Isp(t) = \Isp_0 \cdot \frac{\cstar_{\text{actual}}(t)}{\cstar_{\text{actual,0}}} \cdot \frac{C_f(t)}{C_{f,0}}
\end{equation}

\subsubsection{Chamber Pressure Evolution}

Chamber pressure changes due to geometry:
\begin{equation}
\Pc(t) = \Pc_0 \cdot \sqrt{\frac{A_{\text{throat,0}}}{A_{\text{throat}}(t)}} \cdot \frac{\cstar_{\text{actual}}(t)}{\cstar_{\text{actual,0}}}
\end{equation}

\subsubsection{Mixture Ratio Drift}

Tank pressure depletion affects mixture ratio:
\begin{equation}
\MR(t) = \MR_0 \cdot \sqrt{\frac{P_{\text{tank,O}}(t) - \Pc(t)}{P_{\text{tank,O,0}} - \Pc_0}} \cdot \sqrt{\frac{P_{\text{tank,F,0}} - \Pc_0}{P_{\text{tank,F}}(t) - \Pc(t)}}
\end{equation}

\subsection{Degradation Rates}

Linear degradation approximation:
\begin{equation}
F(t) = F_0 + \dot{F} \cdot t
\end{equation}

where:
\begin{equation}
\dot{F} = \frac{F_{\text{final}} - F_0}{t_{\text{burn}}}
\end{equation}

Degradation rate from linear fit:
\begin{equation}
\dot{F} = \text{slope of linear regression: } F(t) \sim t
\end{equation}

\subsection{Total Impulse}

\begin{equation}
I_{\text{total}} = \int_0^{t_{\text{burn}}} F(t) \, dt
\end{equation}

Numerical integration (trapezoidal rule):
\begin{equation}
I_{\text{total}} \approx \sum_{i=1}^{N-1} \frac{F(t_i) + F(t_{i+1})}{2} \cdot (t_{i+1} - t_i)
\end{equation}

Average thrust:
\begin{equation}
F_{\text{avg}} = \frac{I_{\text{total}}}{t_{\text{burn}}}
\end{equation}

\subsection{Mission Performance}

\subsubsection{Total Delta-V}

From rocket equation:
\begin{equation}
\Delta V = \Isp_{\text{avg}} \cdot g_0 \cdot \ln\left(\frac{m_0}{m_f}\right)
\end{equation}

where $m_0$ and $m_f$ are initial and final masses.

\subsubsection{Mass Flow History}

\begin{equation}
\mdot(t) = \mdot_0 \cdot \sqrt{\frac{\Pc(t) A_{\text{throat}}(t)}{\Pc_0 A_{\text{throat,0}}}} \cdot \frac{\cstar_0}{\cstar(t)}
\end{equation}

Propellant consumed:
\begin{equation}
m_{\text{prop}}(t) = \int_0^t \mdot(\tau) \, d\tau
\end{equation}

\subsubsection{Burnout Time}

Predict time to empty tanks:
\begin{equation}
t_{\text{burnout}} = \frac{m_{\text{prop,total}}}{\mdot_{\text{avg}}}
\end{equation}

or when thrust drops below threshold:
\begin{equation}
F(t_{\text{burnout}}) < F_{\text{threshold}}
\end{equation}

\subsection{Failure Prediction}

\subsubsection{Geometry Limits}

Throat area growth limit:
\begin{equation}
\frac{A_{\text{throat}}(t)}{A_{\text{throat,0}}} > \xi_{\text{threshold}} \quad \Rightarrow \quad \text{Failure}
\end{equation}

Typical threshold: $\xi_{\text{threshold}} = 1.5$--$2.0$ (50--100\% growth).

Expansion ratio limit:
\begin{equation}
\eps(t) < \eps_{\min} \quad \Rightarrow \quad \text{Underexpansion / Failure}
\end{equation}

\subsubsection{Performance Limits}

Thrust drop:
\begin{equation}
\frac{F(t)}{F_0} < 0.5 \quad \Rightarrow \quad \text{Failure}
\end{equation}

Chamber pressure drop:
\begin{equation}
\frac{\Pc(t)}{\Pc_0} < 0.3 \quad \Rightarrow \quad \text{Failure}
\end{equation}

\subsubsection{Recession Limits}

Total recession:
\begin{equation}
r_{\text{total}}(t) = \int_0^t \dot{r}(\tau) \, d\tau > r_{\max} \quad \Rightarrow \quad \text{Burnthrough}
\end{equation}

where $r_{\max}$ is material thickness limit.

\section{Numerical Methods}

\subsection{Root Finding}

\subsubsection{Brent's Method}

For chamber pressure solution:
\begin{enumerate}
    \item Bracket root: $[\Pc_{\min}, \Pc_{\max}]$ where $R(\Pc_{\min}) \cdot R(\Pc_{\max}) < 0$
    \item Check convergence: $|R(\Pc)| < \text{tolerance}$
    \item If not converged:
    \begin{itemize}
        \item Try inverse quadratic interpolation (if safe)
        \item Otherwise use secant method
        \item If neither safe, use bisection
    \end{itemize}
    \item Update bracket and repeat
\end{enumerate}

Bracket validation:
\begin{equation}
R(\Pc_{\min}) \cdot R(\Pc_{\max}) < 0 \quad \text{(required)}
\end{equation}

Convergence check:
\begin{equation}
|R(\Pc)| < \text{tol}_{\text{abs}} \quad \text{or} \quad |R(\Pc)| / |R(\Pc_{\text{initial}})| < \text{tol}_{\text{rel}}
\end{equation}

\subsubsection{Newton's Method}

Alternative method (requires derivative):
\begin{equation}
\Pc_{n+1} = \Pc_n - \frac{R(\Pc_n)}{R'(\Pc_n)}
\end{equation}

with bracketing to ensure $\Pc \in [\Pc_{\min}, \Pc_{\max}]$.

\subsection{Newton-Raphson for Mach Number}

Area-Mach relation:
\begin{equation}
f(M) = \frac{1}{M} \left[\frac{2}{\gamma + 1} \left(1 + \frac{\gamma - 1}{2} M^2\right)\right]^{(\gamma + 1)/(2(\gamma - 1))} - \eps = 0
\end{equation}

Derivative:
\begin{align}
f'(M) &= -\frac{f(M)}{M} + f(M) \cdot \frac{\gamma + 1}{2(\gamma - 1)} \cdot \frac{(\gamma - 1) M}{\gamma + 1 + (\gamma - 1) M^2 / 2} \\
&= -\frac{f(M)}{M} + f(M) \cdot \frac{(\gamma + 1) M}{2 + (\gamma - 1) M^2}
\end{align}

Iteration:
\begin{equation}
M_{n+1} = M_n - \frac{f(M_n)}{f'(M_n)}
\end{equation}

with step damping for stability:
\begin{equation}
M_{n+1} = M_n - \alpha \cdot \frac{f(M_n)}{f'(M_n)}
\end{equation}

where $\alpha \in (0, 1]$ is damping factor (typically 0.5--1.0).

Bounds enforcement:
\begin{equation}
M_{n+1} = \max(1.01, \min(10.0, M_{n+1}))
\end{equation}

Initial guess:
\begin{equation}
M_0 = \begin{cases}
\frac{\eps^{(\gamma - 1)/2}}{\sqrt{(\gamma + 1)/2}} & \eps > 10 \\
1.0 + \sqrt{2(\eps - 1)/(\gamma + 1)} & \eps \leq 10
\end{cases}
\end{equation}

\subsection{Interpolation}

\subsubsection{Bilinear Interpolation}

For 2D grid $(P_c, \MR)$:
\begin{equation}
f(P_c, \MR) = \sum_{i=0}^1 \sum_{j=0}^1 w_{ij} f_{ij}
\end{equation}

Weights:
\begin{align}
w_{00} &= (1-\xi)(1-\eta) \\
w_{01} &= (1-\xi)\eta \\
w_{10} &= \xi(1-\eta) \\
w_{11} &= \xi\eta
\end{align}

where:
\begin{align}
\xi &= \frac{P_c - P_i}{P_{i+1} - P_i} \\
\eta &= \frac{\MR - \MR_j}{\MR_{j+1} - \MR_j}
\end{align}

\subsubsection{Trilinear Interpolation}

For 3D grid $(P_c, \MR, \eps)$:
\begin{equation}
f(P_c, \MR, \eps) = \sum_{i=0}^1 \sum_{j=0}^1 \sum_{k=0}^1 w_{ijk} f_{ijk}
\end{equation}

Weights:
\begin{equation}
w_{ijk} = w_i(\xi) \cdot w_j(\eta) \cdot w_k(\zeta)
\end{equation}

where:
\begin{align}
\xi &= \frac{P_c - P_i}{P_{i+1} - P_i} \\
\eta &= \frac{\MR - \MR_j}{\MR_{j+1} - \MR_j} \\
\zeta &= \frac{\eps - \eps_k}{\eps_{k+1} - \eps_k}
\end{align}

\subsubsection{Robust Interpolation}

Handles boundary cases and NaN values:
\begin{itemize}
    \item Boundary: clip coordinates to grid bounds
    \item NaN in table: use distance-weighted nearest neighbor
    \item Validation: check all weights sum to 1 and are non-negative
\end{itemize}

\subsection{Numerical Stability}

\subsubsection{Safe Division}

Division with zero checking:
\begin{equation}
\text{result} = \begin{cases}
\frac{a}{b} & |b| > \epsilon \\
\text{default} & |b| \leq \epsilon
\end{cases}
\end{equation}

where $\epsilon$ is machine epsilon or physical threshold.

\subsubsection{Safe Square Root}

\begin{equation}
\text{result} = \begin{cases}
\sqrt{x} & x \geq 0 \\
\text{NaN or default} & x < 0
\end{cases}
\end{equation}

\subsubsection{Safe Logarithm}

\begin{equation}
\text{result} = \begin{cases}
\ln(x) & x > 0 \\
\text{NaN or default} & x \leq 0
\end{cases}
\end{equation}

\subsection{Convergence Diagnostics}

Convergence history tracking:
\begin{equation}
\text{history} = \{R(\Pc_0), R(\Pc_1), \ldots, R(\Pc_n)\}
\end{equation}

Check monotonicity:
\begin{equation}
|R(\Pc_n)| < |R(\Pc_{n-1})| \quad \text{(should be true)}
\end{equation}

Check convergence rate:
\begin{equation}
\frac{|R(\Pc_n)|}{|R(\Pc_0)|} < \text{tolerance} \quad \text{after } n \text{ iterations}
\end{equation}

Minimum iterations for validation:
\begin{equation}
n \geq n_{\min} \quad \text{(typically 3)}
\end{equation}

\section{Validation and Constraints}

\subsection{Dimensional Validation}

All equations must be dimensionally consistent. Base SI units: [m, kg, s, K, A, mol, cd].

Example validation:
\begin{align}
[\text{pressure}] &= [\text{kg}/(\text{m} \cdot \text{s}^2)] \\
[\text{force}] &= [\text{kg} \cdot \text{m}/\text{s}^2] \\
[\text{energy}] &= [\text{kg} \cdot \text{m}^2/\text{s}^2]
\end{align}

\subsection{Physical Constraints}

Input validation (errors raised if violated):
\begin{align}
\Pc &> 0, \quad \Pc < 1 \text{ GPa} \\
\Tc &> 0, \quad \Tc < 5000 \text{ K} \\
\gamma &\in [1.01, 2.0] \\
\MR &\in [0.1, 20] \\
\mdot &> 0, \quad \mdot < 1000 \text{ kg/s} \\
\Lstar &> 0, \quad \Lstar < 10 \text{ m} \\
\cstar &\in [500, 3000] \text{ m/s} \\
\Isp &\in [100, 500] \text{ s}
\end{align}

Output validation:
\begin{itemize}
    \item All outputs checked for finiteness: $\text{isfinite}(x)$
    \item All outputs validated against physical bounds
    \item Errors raised (not clamped) if constraints violated
\end{itemize}

\subsection{Conservation Laws}

\subsubsection{Mass Conservation}

\begin{equation}
\mdot_{\text{in}} = \mdot_{\text{out}}
\end{equation}

For steady-state chamber:
\begin{equation}
\mdot_O + \mdot_F = \mdot_{\text{throat}}
\end{equation}

Validation tolerance: $\epsilon_{\text{mass}} < 10^{-6}$ (relative).

\subsubsection{Energy Conservation}

Steady-flow energy equation:
\begin{equation}
h_{\text{in}} + \frac{1}{2} u_{\text{in}}^2 + q = h_{\text{out}} + \frac{1}{2} u_{\text{out}}^2 + w
\end{equation}

For adiabatic nozzle ($q = 0, w = 0$):
\begin{equation}
h_{\text{chamber}} = h_{\text{exit}} + \frac{1}{2} v_{\text{exit}}^2
\end{equation}

Validation tolerance: $\epsilon_{\text{energy}} < 10^{-3}$ (relative, less strict due to approximations).

\subsubsection{Momentum Conservation}

For thrust calculation:
\begin{equation}
F = \mdot v_{\text{exit}} + (\Pexit - P_{\text{ambient}}) A_{\text{exit}}
\end{equation}

Must equal:
\begin{equation}
F = C_f \Pc A_{\text{throat}}
\end{equation}

Validation tolerance: $\epsilon_{\text{thrust}} < 10^{-2}$ (relative).

\subsection{Choked Flow}

Throat must be choked:
\begin{equation}
M_{\text{throat}} = 1
\end{equation}

This requires:
\begin{equation}
\frac{\Pc}{\Pexit} > \left(\frac{\gamma + 1}{2}\right)^{\gamma/(\gamma - 1)}
\end{equation}

For typical $\gamma \approx 1.2$: $\Pc/\Pexit > 1.83$.

\subsection{Thrust Equation Consistency}

Two methods must agree:
\begin{equation}
F_{\text{momentum+pressure}} \approx F_{C_f} = C_f \Pc A_{\text{throat}}
\end{equation}

Relative error:
\begin{equation}
\epsilon_{\text{rel}} = \frac{|F_{\text{momentum+pressure}} - F_{C_f}|}{|F_{C_f}|} < 10^{-2}
\end{equation}

If error too large, use $C_f$ method as fallback.

\subsection{Physics Validation Framework}

\subsubsection{Validation Result Structure}

Each validation returns:
\begin{itemize}
    \item \texttt{passed}: Boolean (true if valid)
    \item \texttt{message}: String description
    \item \texttt{value}: Actual value checked
    \item \texttt{expected\_range}: Expected bounds
    \item \texttt{severity}: "error", "warning", or "info"
\end{itemize}

\subsubsection{Comprehensive Engine State Validation}

Validates entire engine state:
\begin{equation}
\text{validate\_engine\_state}(\Pc, \MR, \mdot, \cstar, \gamma, \Tc, \Isp, F)
\end{equation}

Returns list of validation results (errors first, then warnings, then info).

Critical errors stop execution; warnings logged but continue.

\section{Time-Series Analysis}

\subsection{Pressure Profile Generation}

Time-varying tank pressures:
\begin{equation}
P_{\text{tank}}(t) = P_{\text{tank,0}} \cdot \left(1 - \frac{t}{t_{\text{burn}}}\right)^n
\end{equation}

where $n$ depends on propellant properties and tank geometry.

Linear approximation:
\begin{equation}
P_{\text{tank}}(t) = P_{\text{tank,0}} - \dot{P} \cdot t
\end{equation}

where $\dot{P} = \mdot / (V_{\text{tank}} \rho)$ is pressure drop rate.

\subsection{Array Evaluation}

For time-series simulation:
\begin{equation}
\text{results}(t_i) = \text{runner.evaluate}(P_{\text{tank,O}}(t_i), P_{\text{tank,F}}(t_i))
\end{equation}

for $i = 0, 1, \ldots, N$ time points.

Efficient implementation reuses scalar code path for each time point, guaranteeing identical physics.

\section{Software Implementation}

\subsection{Modular Architecture}

\begin{itemize}
    \item \textbf{Configuration}: Pydantic schemas with validation
    \item \textbf{Thermochemistry}: CEA cache with interpolation
    \item \textbf{Feed System}: Loss models, regenerative cooling
    \item \textbf{Injectors}: Pintle, coaxial, impinging implementations
    \item \textbf{Chamber}: Mass flow balance, root finding
    \item \textbf{Combustion}: Finite-rate chemistry, efficiency models
    \item \textbf{Nozzle}: Shifting equilibrium, thrust calculation
    \item \textbf{Cooling}: Regenerative, film, ablative models
    \item \textbf{Analysis}: Stability, burn analysis, validation
\end{itemize}

\subsection{Data Flow}

\begin{enumerate}
    \item Load configuration from YAML
    \item Initialize CEA cache (build if needed)
    \item Create runner with configuration
    \item For each operating point:
    \begin{enumerate}
        \item Solve chamber pressure (Brent's method)
        \item Calculate injector flows
        \item Get thermochemical properties (CEA)
        \item Calculate reaction progress
        \item Calculate nozzle expansion
        \item Calculate thrust
        \item Evaluate cooling systems
        \item Validate results
    \end{enumerate}
    \item Return diagnostics
\end{enumerate}

\subsection{Configuration Schema}

All parameters defined in YAML with Pydantic validation:
\begin{itemize}
    \item Fluid properties (density, viscosity, etc.)
    \item Geometry (injector, chamber, nozzle)
    \item Feed system (pipes, losses, regen)
    \item Cooling systems (regen, film, ablative)
    \item Combustion (CEA settings, efficiency model)
    \item Solver settings (method, tolerance, bounds)
\end{itemize}

\section{Summary}

This complete mathematical model describes the pintle engine system from propellant tanks through combustion to thrust generation. All equations are physics-based with no empirical reference conditions. The model:

\begin{itemize}
    \item Uses actual Arrhenius kinetics for reaction rates (no reference conditions)
    \item Calculates all time scales from first principles (evaporation, mixing, reaction)
    \item Models shifting equilibrium in nozzle expansion (Damköhler number based)
    \item Accounts for finite-rate chemistry and incomplete combustion
    \item Tracks geometry evolution from ablative cooling (chamber, throat, exit)
    \item Models multiple injector types (pintle, coaxial, impinging)
    \item Includes comprehensive cooling systems (regen, film, ablative, graphite)
    \item Provides stability analysis (chugging, acoustic modes, POGO)
    \item Tracks burn performance and degradation
    \item Validates all physical constraints and conservation laws
    \item Uses robust numerical methods with error handling
\end{itemize}

The system is solved iteratively, with chamber pressure determined by mass flow balance, and all downstream properties calculated from actual engine conditions. No clamping or reference-based scaling—all dynamics are physics-based and validated.

\section{Key Features}

\begin{enumerate}
    \item \textbf{Physics-Based}: No empirical reference conditions—all calculations use actual engine state
    \item \textbf{Finite-Rate Chemistry}: Complete reaction progress modeling with Arrhenius kinetics
    \item \textbf{Shifting Equilibrium}: Nozzle expansion accounts for composition shift
    \item \textbf{Geometry Evolution}: Tracks ablative recession of chamber, throat, and nozzle exit
    \item \textbf{Multiple Injectors}: Supports pintle, coaxial, and impinging designs
    \item \textbf{Comprehensive Cooling}: Regenerative, film, and ablative with detailed heat transfer
    \item \textbf{Stability Analysis}: Combustion and feed system stability evaluation
    \item \textbf{Burn Analysis}: Performance degradation, mission metrics, failure prediction
    \item \textbf{Numerical Robustness}: Safe operations, validation, convergence checking
    \item \textbf{Modular Design}: Extensible architecture for additional injectors/cooling methods
\end{enumerate}

\end{document}

